\subsection{Queries on Data Streams}
\label{ss:queries}

Given a data stream formed by user data at each timestamp, i.e., $\CD =
(D_1,D_2,\ldots)$, other than the histogram/frequency query extensively studied
in previous work~\cite{Kellaris2014, Ren2022}(i.e., the query $\bc_t/n_t$ defined in \cref{ss:data_stream}), we
define several novel streaming query tasks.

\begin{definition}[Counting Query]
  Let $c_t(v)$ denote the number of users reporting the private value $v\in\Omega$ at
  time $t$, i.e., $c_t(v)\triangleq |\{u_i\in U\colon v_{it}=v\}|$.
  Given a value $v\in\Omega$ and a positive integer $\Delta$, a counting query
  aims to report the total number of times users report the private value $v$ in
  the most recent $\Delta$ timestamps.
  Formally,
  \[
    Q_t^c(v,\Delta)\triangleq \sum_{t'=t-\Delta +1}^t c_{t'}(v).
  \]
\end{definition}

The counting query is a generalization of the histogram query and has widespread
applications.
Consider the following example.

\begin{example}\label{exam:vehicle}
  Consider a vehicle dispatch management system in a city.
  At each time $t$, the system needs to publish the total number of times
  vehicles passed by a location $v$ over the past ten minutes.
  Assuming a time granularity of one minute, then the query corresponds to
  $Q_t^c(v, 10)$.
\end{example}

As a special case, when $\Delta=1$, $Q_t^c(v,1)=c_t(v)$ which is the number of
users reporting $v$ at time $t$, and $Q_t^c(v,1)/n_t$ becomes the frequency of
value $v$ at time $t$.

\begin{definition}[Range Query]
  Given a value range $[v_1, v_2]$ (where $v_1\leq v_2$) and a positive integer
  $\Delta$, a range query aims to report at each time $t$ the total number of
  times users report private values in $[v_1, v_2]$ in the most recent
  $\Delta$ timestamps.
  Formally,
  \[
    Q_t^r([v_1, v_2],\Delta)\triangleq
    \sum_{t'=t-\Delta +1}^t\sum_{v\in[v_1,v_2]} c_{t'}(v)
  \]
\end{definition}

The range query further generalizes the counting query by allowing the user value to
be in a specified range.
Consider the following example.

\begin{example}\label{exam:disease}
  A disease monitoring center needs to publish daily the total number of
  infected people with age above $60$ over the past week during the COVID
  pandemic.
  Assuming a time granularity of one day, then the query corresponds to $Q_t^r([60,
  \infty), 7)$.
\end{example}

In many real-world scenarios, the server also wants to identify specific
abnormal events in the stream through a class of queries known as event
monitoring.

\begin{definition}[Event Monitoring]\label{def:event_monitoring}
  Let $\alpha$ and $\beta$ be two functions, where $\alpha$ could be a counting
  query or a range query on the stream, and $\beta$ is a Boolean function defined
  on the output of $\alpha$.
  An event monitoring query, denoted by $Q_t^e(\alpha,\beta)$, aims to report at
  each time $t$ a binary result where $1$ denotes the monitored event occurs,
  and $0$ otherwise.
\end{definition}

We illustrate the usefulness of this concept by the following example.
Let $\Delta>0$ denote a specified monitoring period, and we are interested in
monitoring whether counts of a value $v$ significantly increase in two
consecutive monitoring periods.
Then the two functions $\alpha$ and $\beta$ can be formally defined as
\begin{align*}
  \alpha(v,\Delta) &= Q_t^c(v, \Delta) - Q_{t-w+1}^c(v, \Delta), \\
  \beta(x,\vartheta) &= \indr{x>\vartheta},
\end{align*}
where $\indr{B}$ denotes an indicator function, and it returns $1$ if the
condition $B$ is true, otherwise $0$.
The event monitoring query is a compound function on the stream, i.e.,
$Q_t^e(\alpha,\beta)= \beta\circ\alpha$, which reports $1$ if there is a sudden
increase of counts of value $v$ in two consecutive monitoring periods.
It is also straightforward to define $\alpha$ using a range query to monitor
multiple user values.
Therefore, we can use the defined event monitoring queries to monitor the sudden
change in counts in \cref{exam:vehicle,exam:disease}.

\paragraph{Remarks}
% \remarks.
The streaming queries proposed above are generalizations of existing
works~\cite{Kellaris2014,Chen2017} with the emphasis on data streams and aiming
to achieve $w$-event LDP.
Furthermore, for parameter $\Delta$ in these queries, we require $\Delta \leq w$
in order to satisfy $w$-event LDP.
